\clearpage
\fancyhead[L]{\small GAUSS AN VON LINDENAU.}
\SourceLeaf{310}\begin{center}
GAUSS AN VON LINDENAU.
\end{center}

\hrule
\begin{center}
\fontsize{10}{12}\selectfont Monatliche Correspondenz. Band XVIII. S. 173..175. August 1808.
\end{center}
\hrule

\begin{flushright}
Göttingen, den 6. August 1808.
\end{flushright}

\par\indent Seit meinem letzten Brief habe ich die Beobachtungen der neuen Planeten fortgesetzt, sooft es die Umstände erlaubten. Ich freue mich, Ihnen hier mitzuteilen, was ich bisher ermittelt habe. \emph{Juno}, \emph{Vesta} und \emph{Pallas} ließen sich bisher nur mit dem Kreismikrometer beobachten. Die Beobachtungen der Vesta habe ich aber seit dem 13. Juli nicht fortgesetzt, da sie sich von jetzt an schon im Meridian beobachten lässt. Professor \textsc{Harding} ist auch eine Beobachtung der \emph{Ceres} am Quadranten gelungen. Ich werde sie Ihnen später mitteilen [handschriftliche Bemerkung: Juli 25. $13^{\text{u}}\,14^{\text{m}}\,29^{\text{s}}$ \quad $322^\circ\,8'\,57''6$ \quad $27^\circ\,52'\,11''5$]

\par\indent \emph{Pallas} steht jetzt leider die ganze Nacht für unseren Beobachtungsort zu hoch. Ich konnte nur einige Beobachtungen machen, deren Deklinationen aber unsicher sind.

\par\indent Die Beobachtung der \emph{Vesta} vom 2. Juli und die der \emph{Pallas} vom 15. Juli stammen von Herrn \textsc{Harding}, alle übrigen von mir. Der Fehler der Ephemeride beträgt bei \emph{Pallas} $-4'$ in der Rektaszension. Die Deklination stimmt auf $1'$ genau. Bei \emph{Juno} ist die Differenz in der Rektaszension inzwischen auf $13'$ angewachsen. Die vier letzten Beobachtungen der \emph{Juno} halte ich alle für sehr gut, sofern die Vergleichssterne in der Hist. Cél. gut beobachtet worden sind.

\Needspace{13\baselineskip}
\begin{center}
\emph{Beobachtungen der Pallas.}
\end{center}


\begin{center}\fontsize{10}{12}\selectfont\setlength{\tabcolsep}{3pt}\renewcommand{\arraystretch}{1.15}\begin{tabular}{lrrr}
1808 & \shortstack{Mittlere Zeit\\in Göttingen} & \shortstack{Scheinbare\\Rektaszension} & \shortstack{Scheinbare\\Deklination}\\\hline
Juli 13 & $12^{\text{u}}\;3^{\text{m}}\;17^{\text{s}}$ & $300^\circ\;21'\;49''$ & $18^\circ\;41'\;44''$ N.\\
15 & $12\;1\;28$ & $299\;57\;3$ & $18\;32\;40$\\
25 & $10\;9\;44$ & $297\;57\;37$ & $17\;40\;10$\\
\end{tabular}\end{center}


\Needspace{13\baselineskip}
\begin{center}
\emph{Beobachtungen der Juno.}
\end{center}


\begin{center}\fontsize{10}{12}\selectfont\setlength{\tabcolsep}{3pt}\renewcommand{\arraystretch}{1.15}\begin{tabular}{lrrr}
1808 & \shortstack{Mittlere Zeit\\in Göttingen} & \shortstack{Scheinbare\\Rektaszension} & \shortstack{Scheinbare\\Deklination}\\\hline
Juni 20 & $11^{\text{u}}\;49^{\text{m}}\;0^{\text{s}}$ & $315^\circ\;29'\;34''$ & $2^\circ\;16'\;23''$ S.\\
22 & $12\;0\;45$ & $315\;23\;1$ & $2\;14\;29$\\
Juli 6 & $12\;42\;23$ & $314\;1\;0$ & $2\;19\;7$\\
30 & $10\;51\;17$ & $309\;38\;9$ & $3\;57\;13$\\
31 & $11\;40\;57$ & $309\;24\;53$ & $4\;3\;48$\\
Aug. 4 & $10\;42\;53$ & $308\;33\;44$ & $4\;30\;32$\\
5 & $10\;53\;46$ & $308\;20\;39$ & $4\;37\;36$\\
\end{tabular}\end{center}


\par\indent [Im Original sind auch Beobachtungen der \emph{Vesta} vom 22. Juni sowie vom 1., 2. und 13. Juli beigefügt. Sie finden sich auch in der weiter unten auf Seite 309 abgedruckten Mitteilung der Gött. gel. Anz. vom 10. Nov. 1808.]

\par\indent Ich habe bereits levi calamo aus meinen Beobachtungen die Opposition der \emph{Juno} berechnet und die Elemente vorläufig verbessert. Eine kleine Änderung genügt, um alle vier bisher beobachteten Oppositionen mit einer reinen Ellipse sehr genau darzustellen. Sie besteht hauptsächlich in einer leichten Verringerung der Länge des Perihels und einer leichten Erhöhung der mittleren Bewegung. Bald mehr darüber.

\par\indent Unsere Sternwarte hat vom König mehrere astronomische Instrumente geschenkt bekommen. Ein gregorianisches Teleskop von \textsc{Chevalier} und ein Marinechronometer von \textsc{Louis Berthoud} sind schon an\SourceLeaf{311}gekommen. Soweit sich aus den bisherigen Erfahrungen schließen lässt, hat das Chronometer einen ziemlich gleichmäßigen Gang. Freilich wurde sein Gang bisher nur bei beständiger Ruhe beobachtet. Bald werden wir versuchen, damit eine geografische Länge zu bestimmen. [Diese Mitteilung über die Instrumente ist auch in den Gött. gel. Anzeigen, Stück 127, Seite 1265, vom 8. August 1808 mit dem Zusatz abgedruckt „Die übrigen Instrumente, die bisher in Paris waren, dürfen wir ebenfalls bald erwarten.“ \textsc{Schering}.]

